\documentclass{article}
\usepackage{preamble}

\begin{document}
\header{Filemon Mateus}{CS 6140}{Document Similarity \& Hashing}{Homework \# 2}{\today}
\begin{enumerate}[leftmargin={*}, font={\bf}, label={\arabic*.}, ref={\arabic*}]
  \item \label{qst:1}
    \begin{enumerate}[label={(\Alph*)}, ref={\theenumi(\Alph*)}]
      \item \label{qst:1A}
        Results from distinct $n$-Grams on each review:
        \begin{table}[H]
          \centering
          \begin{tabular}{|ccccc|}
            \hline
            {\tt n-Grams} & {\tt D1} & {\tt D2} & {\tt D3} & {\tt D4} \\ \hline \hline
            2-grams (char) & 333 & 461 & 401 & 154 \\ \hline
            3-grams (char) & 933 & 1714 & 1428 & 256 \\ \hline
            2-grams (word) & 330 & 875 & 571 & 54 \\ \hline
          \end{tabular}
          \caption{
            Distinct $n$-Grams for each review. See \autoref{lst:build-ngrams} for \\
            implementation details.
          }
          \label{tbl:build-ngram-table}
        \end{table}
        The results from \autoref{tbl:build-ngram-table} were obtained by creating a
        {\tt CountVectorizer} object with the following parameters, then applying the
        {\tt fit} and {\tt get\_feature\_names\_out} methods:

        \begin{lstlisting}
-- CountVectorizer(analyzer = 'char', ngram_range = (2,2))
-- CountVectorizer(analyzer = 'char', ngram_range = (3,3))
-- CountVectorizer(analyzer = 'word', ngram_range = (2,2))
        \end{lstlisting}

      \item \label{qst:1B}
        A subroutine for generalized set similarity is implemented below:
        \begin{lstlisting}[
          label={lst:generalized-set-similarity-measure},
          caption={Generalized set similarity measure.}
        ]
def generalized_set_similarity_measure(A, B, C, D, x, y, z, z_prime):
    a = len(A.intersection(B))
    b = len(A.union(B).union(C).union(D)) - len(A.union(B))
    c = len(A.symmetric_difference(B))
    d = x * a + y * b + z * c
    e = x * a + y * b + z_prime * c
    return d / e
        \end{lstlisting}

        \bigskip

        Below are the Jaccard similarities corresponding to every pair:
        \vspace{-5\parskip}

        \begin{minipage}{\linewidth}
          \begin{table}[H]
            \centering
            \begin{tabular}{|ccccccc|}
              \hline
              {\tt n-Grams} & {\tt JS(D1,D2)} & {\tt JS(D1,D3)} & {\tt JS(D1,D4)} & {\tt JS(D2,D3)} & {\tt JS(D2,D4)} & {\tt JS(D3,D4)} \\ \hline \hline
              2-grams (char) & 0.59118 & 0.62031 & 0.40751 & 0.66409 & 0.30851 & 0.34709 \\ \hline
              3-grams (char) & 0.32682 & 0.31021 & 0.17839 & 0.40080 & 0.11932 & 0.12567 \\ \hline
              2-grams (word) & 0.01516 & 0.01923 & 0.00524 & 0.02263 & 0.00324 & 0.00160 \\ \hline
            \end{tabular}
            \caption{
              Jaccard similarity on every document pair. See \autoref{lst:build-ngrams} \\
              for implementation details.
            }
          \end{table}
        \end{minipage}
    \end{enumerate}

  \item \label{qst:2}
    \begin{enumerate}[label={(\Alph*)}, ref={\theenumi(\Alph*)}]
      \item \label{qst:2A}
        Results of the approximation of Jaccard similarity base on {\tt fast-min-hash}
        algorithm:
        \begin{table}[H]
          \centering
          \begin{tabular}{|cc|}
            \hline
            {\tt t} & {\tt $\sim$JS(D1,D2)} \\ \hline \hline
            20 & 0.85000 \\ \hline
            60 & 0.90000 \\ \hline
            150 & 0.92000 \\ \hline
            300 & 0.90000 \\ \hline
            600 & 0.89833 \\ \hline
          \end{tabular}
          \caption{
            Approximation of Jaccard similarity based on {\tt fast-min-hash} algorithm. \\
            See \autoref{lst:fast-min-hash} for implementation details.
          }
        \end{table}

      \item \label{qst:2B}
        Benchmark results:
        \begin{figure}[H]
          \centering
          \begin{tikzpicture}
            \begin{axis}[
              title={Benchmark Results for MinHash Algorithm},
              xlabel={$t$-values},
            ]
              \addplot[
                thin,
                color=red
              ] table[col sep=comma] {dev-dir/week2/data/error.csv};
              \addlegendentry{Residual Error};
              \addplot[
                thin,
                color=green
              ] table[col sep=comma] {dev-dir/week2/data/times.csv};
              \addlegendentry{Running Time};
            \end{axis}
          \end{tikzpicture}
          \caption{
            Benchmark results from {\tt fast-min-hash}. Plotted values correspond to \\
            expectation on $100$ trials. See \autoref{lst:fast-min-hash} for implementation
            details.
          }
          \label{fig:fast-min-hash-benchmark}
        \end{figure}

        From \autoref{fig:fast-min-hash-benchmark}, we see that $t \approx 35$ produces the
        best approximation with the least runtime overhead. In general, a choice $t = 35$
        would be appropriate if dealing with massively large text files. For the datasets
        provided, however, choosing $t \in [300, 400]$ would be ideal as values past $t >
        400$ results in diminishing improvements in accuracy while continuing to introduce
        significant computational overhead to the running time.
    \end{enumerate}
\end{enumerate}

\newpage

\section*{Appendix}

\lstinputlisting
  [caption={\sf build-ngrams.py}, label={lst:build-ngrams}]
  {dev-dir/week2/build-ngrams.py}

\lstinputlisting
  [caption={\sf fast-min-hash.py}, label={lst:fast-min-hash}]
  {dev-dir/week2/fast-min-hash.py}
\end{document}
